\section{Introduction}

When a computer shows that no decision rule can meet all the requirements we
asked for, the next question is not simply whether the computation was
correct. We also need to know what the result allows us to say. Which
requirement is responsible for the conflict? Which change would actually
repair it? Which parts of the answer are backed by the recorded evidence, and
which parts depend on modeling choices made before the check was run? This
paper studies that question in a small social-choice case and argues that a
formal artifact should report only the claims that its evidence supports.

This matters because formal artifacts increasingly sit between technical
evidence and downstream decisions. In formal social choice, they may report
that a collection of requirements is impossible or that a proposed repair is
available. In software verification, they may separate a real inconsistency
from a modeling artifact. In AI-mediated governance workflows, they may help
explain why a rule, policy, or constraint set failed. In all of these settings,
the problem is not only whether an internal check succeeded. The further
question is what public or scientific claim that check licenses. A correct
artifact can still be over-reported if its certificate, encoding assumptions,
and prose interpretation are treated as though they were the same object.

The running example is Sen24, a deliberately small instance from social choice
theory with two voters and four alternatives. The case is small enough for the
relevant finite evidence to be audited completely inside a declared finite
interface, but rich enough to expose the gap between an impossibility result
and a justified repair report. It contains multiple requirements, multiple
possible repairs, and nontrivial ways to group those repairs for explanation.
M4 uses this case because it is a controlled setting in which the evidence
boundary can be made explicit. It
does not claim to solve social choice in general, and it does not present the
finite case as a substitute for a general theorem.

The main contribution is methodological. M4 develops a claim-boundary audit for
finite social-choice artifacts: it separates finite evidence, declared semantic
assumptions, partial formal proof, and future correctness obligations. The
point is not merely to record that an artifact was checked. The point is to say
which repair and reporting claims are certified under the declared scope,
which semantic targets are backed by formal definitions, and which stronger
correctness claims remain outside the present paper.

The technical instantiation is the declared M4 Sen24 encoding. Under that
encoding, the finite audit is organized by bundled masks and obstruction
shapes. It records the 48-cell audit, distinguishes SAT-side repair
classification from UNSAT-side orbit/fiber exactness, and includes support
truncation for shape-blind reporting. These are paper-facing certificate
claims under the declared encoding, not general structural theorems. The paper
also uses a partial Lean-backed semantic target for the central right
condition, \code{RightAtomSemantics}, which packages endpoint distinctness
with Lean \code{Decisive}. The implemented bridge supports orientation
invariance and the sufficient direction from two fixed right conditions with
distinct voters to Lean \code{MINLIB}. It does not prove that the Python or
generated artifact layers encode that target.

This paper is deliberately not an end-to-end verification result. The general
Sen theorem belongs to the M2 semantic bridge and is not a new M4 claim. The
M4 claim is narrower and more explicit: under the declared finite encoding,
the paper audits which repair and reporting statements are licensed, while
marking the remaining artifact-correctness bridges as future correctness
obligations. Section~\ref{sec:limitations} records the full non-claim list.

The contributions are:

\begin{itemize}[leftmargin=*]
\item a claim-boundary framing for finite social-choice artifacts, focused on
  what repair and reporting claims the evidence licenses;
\item a complete finite audit under the declared M4 Sen24 encoding, stated as
  a paper-facing certificate rather than as an unrestricted theorem;
\item a paper-facing organization of the repair and reporting structure that
  separates the finite evidence layer from semantic assumptions;
\item a Lean-backed semantic target for the central right condition, including
  the sufficient bridge to \code{MINLIB} for two fixed conditions with distinct
  voters;
\item an explicit separation of crossed bridges from future correctness
  obligations, including the uncrossed Python/CNF and artifact-level semantic
  bridges.
\end{itemize}

Section~2 situates the case in the M1--M4 artifact line and identifies which
earlier results are background rather than M4 contributions. Section~3 states
the declared encoding and scope. Section~4 presents the finite audit as a
paper-facing certificate. Section~5 describes the Lean-backed semantic target
and its limits. Section~6 collects the explicit claim boundaries and bridge
statuses. Section~7 discusses the broader value of auditable abstraction
contracts. Section~8 positions the work against related literatures.
Section~9 records limitations and future correctness obligations. Appendix~A
records reproducibility and artifact-binding information.
